\section{Experimental protocol}
\label{sec:protocol}

This section defines the data split in time, the users, the selection rule, the statistical tests, the leakage controls and
the baselines. The numbers of the results are in Section~\ref{sec:results}. Every number in this section is either a count
from a data file or a configuration count from a run log, and each is given with its source.

\subsection{Temporal protocol and the role of the random-exposure data}
\label{sec:protocol-temporal}

The standard log is the policy log of the platform. The random-exposure log is the set of impressions in which a feed slot was
replaced by a video drawn uniformly from the pool $\mathcal I$. The random-exposure
period starts on 2022-04-22 at 00:00 UTC+8, and this instant is the cut-off $T_0$ of the protocol
(Section~\ref{sec:method-notation}).

Every quantity used for training or for selection is computed before $T_0$. This includes the in-pool click matrix and the
standard log used by the backbones, the lifelong tables of Section~\ref{sec:method-leaf}, the DAX fit, the popularity offset,
and the negative sampler. The training log is the standard log before the cut-off, and the test lists are the random
impressions after it. A pair that was shown before $T_0$ is excluded from its user's candidate list, so a candidate is never
a pair that the backbone saw as a positive or as an impression.

The reason for the cut-off is the feedback path. After a random impression the policy keeps logging, so a standard log that
extends past $T_0$ can contain the user's later responses to the randomly inserted video. A protocol that trains on that log
can therefore learn from the test outcome. The concurrent protocol, which uses the whole standard log, is reported only as a
robustness row in Section~\ref{sec:results}.

The random-exposure data are used for two purposes only. They provide the candidate lists and labels of the validation users,
which select the configuration, and of the test users, which measure the score. No model reads the random log for training.
This is the setting that matches the uniform-insertion target. A method that trained on uniform labels would be measured on
its own training target, and the comparison would no longer test the use of the policy log. Validation labels enter only
through the selection rule, and the test labels only through the reported test value.
% source: drup/data2.py load_kuairand_rich (train_until, candidate construction, meta leak checks); drup/leaf.py check_temporal_protocol

\subsection{Datasets}
\label{sec:protocol-data}

The primary dataset is KuaiRand-27K \citep{gao2022kuairand}. Its users are all users that appear in both the standard log
and the random log, and its pool is the set of random-log videos. The secondary dataset is KuaiRand-1K, from the same release
with 1,000 users, 993 of whom appear in both logs. Its user identifiers are the same as those of the 27K release, so the 1K results are a replication at a
smaller scale and not an independent sample. Table~\ref{tab:data} gives the sizes under the temporal protocol.

\begin{table}[t]
\centering
\caption{Sizes of the two KuaiRand settings under the temporal protocol. A cell marked TBD is not in a result file.}
\label{tab:data}
\footnotesize
\begin{tabular}{p{0.36\textwidth}p{0.15\textwidth}p{0.15\textwidth}p{0.32\textwidth}}
\toprule
Quantity & KuaiRand-27K & KuaiRand-1K & Source \\
\midrule
Users in both logs (dataset users) & 27,077 & 993 & run logs (27K: \texttt{dax\_\allowbreak{}lgn\_\allowbreak{}post.log}; 1K: \texttt{din\_\allowbreak{}1k.log}) \\
Pool items & 7,583 & 7,388 & same \\
Pre-cut training clicks & 523,176 & 18,873 & \texttt{din\_\allowbreak{}27k.log}; \texttt{din\_\allowbreak{}1k.log} \\
Pre-cut training impressions & \{TBD:temporal impressions 27K\} & \{TBD:temporal impressions 1K\} & not in a file \\
Users with a pre-cut lifelong history & 26,723 & 977 & \texttt{din\_\allowbreak{}27k.log}; \texttt{din\_\allowbreak{}1k.log} \\
Out-of-pool impressions before $T_0$ & 135,219,168 & \{TBD:1K pre-cut out-of-pool\} & \texttt{lifelong\_\allowbreak{}27k\_\allowbreak{}summary.txt}, all users of the release \\
Random-exposure rows & 1,186,059 & 42,944 & \texttt{link\_\allowbreak{}users.txt}; \texttt{stats\_\allowbreak{}1k.txt} (993 users) \\
Test users (candidate lists) & 27,077 & 993 & run logs \\
Test users with a click, pooled & 22,612 & 842 & \texttt{stageB.txt} Section 2; \texttt{stats\_\allowbreak{}1k.txt} \\
Test users with a click, per split & 15,809 to 15,913 & \{TBD:1K per split\} & \texttt{stageB.txt} Section 2 \\
Candidates per user (temporal) & \{TBD:27K candidates per user\} & \{TBD:1K candidates per user\} & not in a file \\
\bottomrule
\end{tabular}
\vspace{2pt}

\raggedright\scriptsize
Sources are in \texttt{results/} (\texttt{v9\_probes}, \texttt{v11}, \texttt{logs\_v11}, \texttt{logs\_v12}). The 42,944 rows are the
random-log impressions of the 993 users, and the random log does not depend on the protocol.
\end{table}

Coat \citep{schnabel2016} and Yahoo! R3 \citep{marlin2009collaborative} contain ratings that users gave, together with a
randomly chosen set of ratings for the test part. They record no impressions that were shown and left unrated, and they have no
history on other items. DAX needs impressions with and without a click in each scenario, and LEAF needs a user's history on
videos outside the test pool. Neither component can be run on these data, and running the backbones alone would test a
different question. KuaiRec \citep{gao2022kuairec} has a near-complete small user-video matrix and a larger sparse log of watch
ratios. It has no scenario field among its interaction fields, and the user's history is confined to its own matrix, so the
history on videos outside the test pool is not available. Its daily item counts are platform aggregates and not user
impressions. The Open Bandit Dataset has no persistent user identifiers, so no lifelong history can be formed for it.

The motivation for a randomised test is the headroom of the classic benchmarks. A leave-one-user-out item oracle uses test labels,
so it is a diagnostic and not a method. On KuaiRec, the earlier diagnostics of the item oracle relative to the best training-free
prior are not yet in a result file: \{TBD:results/v12/saturation\_kuairec.txt oracle nDCG@20, prior nDCG@20, best gain\}. On
KuaiRand-27K the same kind of diagnostic is in a file, and it is over all users: the item oracle reaches nDCG@10 0.5072, the log
prior 0.3884 and the random ranking 0.3329. A rank-64 model on the residuals of the uniform item means moves the item-mean value
from 0.5850 to 0.5851 (\texttt{v9\_probes/headroom\_kuairand\_pure.txt}). The residual headroom therefore lies at the item level, and
that is the level at which DAX operates.

\subsection{Splits, users and per-split selection}
\label{sec:protocol-splits}

Each split partitions the users of the test list into validation and test users, by user and not by impression. The split with
index $s\in\{0,\dots,4\}$ takes a random permutation of the users with seed $s$, and puts the first $\operatorname{round}(0.3\,m)$
users into validation. For KuaiRand-27K this gives $\operatorname{round}(0.3\times27{,}077)=8{,}123$ validation users and 18,954
test users in each split. For KuaiRand-1K the first split has 298 validation and 695 test users
(\texttt{v9\_\allowbreak{}probes/stats\_\allowbreak{}1k.txt}). The two parts share no user, and the code asserts this for every split
(\texttt{drup/leaf.py protocol\_\allowbreak{}split}). Users with no click among their random impressions have no defined nDCG and are not
evaluated, so the evaluated test users per split are the numbers in Table~\ref{tab:data}.

The training runs do not depend on the split. Each backbone is trained once for each hyperparameter setting on the whole
pre-cut log, and every split reads the same training trajectory. The splits change only which users select the epoch, the
configuration and the weight $\gamma$, and which users are scored. This is why the split-level test of
Section~\ref{sec:protocol-stats} is needed.

For each split, the selected configuration is the one with the largest validation nDCG@10. The configuration contains the
backbone setting, the epoch, the DAX configuration and $\gamma$ where they apply. The test labels are read only for the
configuration that is the best on validation so far, and the value stored for the split is the value of the final best
configuration. No test label influences the choice of any configuration.

The configuration count of a method, reported as $n_{\mathrm{configs}}$ in the run logs, is the number of backbone trainings times
the number of DAX configurations times the number of $\gamma$ values. For a backbone without DAX it is the number of trainings
or hyperparameter settings. Epochs are selected inside each training but are not counted as configurations. They are counted
in the number of evaluations $n_{\mathrm{evals}}$ instead. For example, LEAF-LGN+DAX has $12\times32\times10=3{,}840$ configurations
and 86,400 evaluations, and QP family+DAX has $10\times32\times10=3{,}200$ configurations. Table~\ref{tab:baselines} gives these
counts for every method, and the search budgets differ. The fair baseline is the one with the same DAX module.

\subsection{Statistical tests}
\label{sec:protocol-stats}

\paragraph{Paired per-user differences.} For two methods $A$ and $B$, a user's difference is the mean, over the splits in
which the user is evaluated, of the difference of the two nDCG@10 values. The reported quantity is the mean of these user
differences, multiplied by $10^{3}$, with a 95\% half-width from the two-sided t interval over users. The p-value is that of a
two-sided one-sample t test on the user differences. Users appear in several test parts and share one training run, so the
user-level t test is approximate, and it is reported together with the split-level test.

\paragraph{Split-level test.} The Nadeau-Bengio corrected resampled t test \citep{nadeau2003} is applied to the five split means
$\bar d_s$ of the difference $d_s=\mathrm{nDCG}_A(s)-\mathrm{nDCG}_B(s)$. With $J=5$ splits, the variance of the mean is
estimated as
\begin{equation}
  \widehat{\mathrm{Var}}=s_d^2\Bigl(\frac1J+\frac{n_{\mathrm{test}}}{n_{\mathrm{val}}}\Bigr)
  =s_d^2\Bigl(\frac15+\frac{7}{3}\Bigr),
\end{equation}
where $s_d^2$ is the sample variance of the $J$ split differences and the ratio is $0.7/0.3$. The statistic is
$t=\bar d/\sqrt{\widehat{\mathrm{Var}}}$, compared with a t distribution with $J-1=4$ degrees of freedom. We take the validation part
as the training part of the original correction and the test part as its test part. The ratio $7/3$ is larger than the
alternative $3/7$, so this choice inflates the variance more and is the more conservative of the two. With these settings, the
two-sided test at level 0.05 needs a mean-to-standard-deviation ratio of the split differences of at least
$t_{0.975,4}\sqrt{1/5+7/3}\approx4.42$.

\paragraph{Holm adjustment.} The p-values of a family are adjusted by the step-down procedure of Holm \{TBD:citation key holm1979,
not in the allowed bibliographies\}. Each table names its family. The primary family for gate B1 has size five and was fixed before the
Stage B runs: LEAF-LGN+DAX against the QP family, the QP family with DAX, the post-hoc LEAF-LGN backbone, r-AdjNorm and LightGCN. The
table family of the deep-baseline extension contains every comparator that is present in the table, and it was defined after the
Stage B results. It has seven members with the deep baselines. The Stage 1 family for the closed-form comparison contains the two LEAF
variants against four baselines, so it has eight members. The same adjustment is applied to the t p-values and to the
Nadeau-Bengio p-values, separately.

\paragraph{Gate and reading.} The gate B1 rule was written in \texttt{paper/v7/DESIGN.md} (Section 7, dated 2026-11-08) before the
Stage B runs. It requires that LEAF-LGN+DAX exceed the QP family by at least $15\times10^{-3}$ per user, and that both Holm-adjusted
p-values be below 0.05. The comparison with the fair baseline, QP family with DAX, has no pre-set margin. It is read with the same
significance level. We report the t and Nadeau-Bengio p-values only. No Wilcoxon test is used.
% source: experiments/significance.py (nadeau_bengio, holm, user_diffs); experiments/final_v11.py (peru, compare, pair, B1 family)

\subsection{Leakage checks and controls}
\label{sec:protocol-controls}

Table~\ref{tab:controls} lists every test in \texttt{tests/} and the checks printed by the runs and the controls of the
analysis. The unit tests were run on 2026-11-13 and all 15 of them passed (6 in \texttt{tests/test\_\allowbreak{}lifelong.py} and 9 in
\texttt{tests/test\_\allowbreak{}dax.py}). The kind of each number is given in the table. A ``main'' number is a validation-selected result of
the main setting, a ``probe'' number is a post-hoc value on stored scores, a ``diagnostic'' number uses random-exposure labels of
the validation users and is never a method or a selection input, and a ``code'' or ``unit'' row states what a test or a code path
asserts.

\begin{table}[p]
\centering
\caption{Leakage checks and controls, one row per unit test and per check.}
\label{tab:controls}
\fontsize{7.2}{8.6}\selectfont
\renewcommand{\arraystretch}{0.8}
\setlength{\tabcolsep}{2pt}
\begin{adjustbox}{max width=\textwidth}
\begin{tabular}{p{0.41\textwidth}p{0.57\textwidth}}
\toprule
Check or test & What it asserts or shows; kind and source in brackets \\
\midrule
\multicolumn{2}{l}{\emph{Unit tests in \texttt{tests/test\_\allowbreak{}lifelong.py} (6 tests, all passed on 2026-11-13)}} \\
\texttt{test\_\allowbreak{}builder\_\allowbreak{}uses\_\allowbreak{}out\_\allowbreak{}of\_\allowbreak{}pool\_\allowbreak{}history\_\allowbreak{}strictly\_\allowbreak{}before\_\allowbreak{}cut} & Out-of-pool events before the cut-off only; synthetic user: 5 impressions, 2 clicks \emph{[unit; builder]} \\
\texttt{test\_\allowbreak{}pre\_\allowbreak{}window\_\allowbreak{}tables\_\allowbreak{}ignore\_\allowbreak{}post\_\allowbreak{}cut\_\allowbreak{}and\_\allowbreak{}pool\_\allowbreak{}rows} & Pre-cut tables unchanged when post-cut and pool rows are altered \emph{[unit; builder]} \\
\texttt{test\_\allowbreak{}tables\_\allowbreak{}ignore\_\allowbreak{}npz\_\allowbreak{}rows\_\allowbreak{}of\_\allowbreak{}users\_\allowbreak{}outside\_\allowbreak{}the\_\allowbreak{}dataset} & Rows of an unknown user identifier leave the tables unchanged \emph{[unit; \texttt{leaf.py}]} \\
\texttt{test\_\allowbreak{}lift\_\allowbreak{}and\_\allowbreak{}familiarity\_\allowbreak{}match\_\allowbreak{}the\_\allowbreak{}formula} & Tables equal $(c+5r)/(n+5)-r$ and $\log(1+n)$, up to scale \emph{[unit; \texttt{leaf.py}]} \\
\texttt{test\_\allowbreak{}labels\_\allowbreak{}of\_\allowbreak{}test\_\allowbreak{}and\_\allowbreak{}validation\_\allowbreak{}never\_\allowbreak{}enter\_\allowbreak{}the\_\allowbreak{}channels} & No label argument; flipping validation and test labels changes nDCG only \emph{[unit; \texttt{leaf.py}]} \\
\texttt{test\_\allowbreak{}shuffled\_\allowbreak{}users\_\allowbreak{}remove\_\allowbreak{}the\_\allowbreak{}lifelong\_\allowbreak{}gain} & Synthetic gain above 0.10 falls below 0.03 in absolute value and below a quarter after shuffling users \emph{[unit; \texttt{leaf.py}]} \\
\multicolumn{2}{l}{\emph{Unit tests in \texttt{tests/test\_\allowbreak{}dax.py} (9 tests, all passed on 2026-11-13)}} \\
\texttt{test\_\allowbreak{}em\_\allowbreak{}likelihood\_\allowbreak{}is\_\allowbreak{}monotone\_\allowbreak{}without\_\allowbreak{}priors} & Observed-data likelihood does not decrease over 25 iterations (negligible priors), three seeds \emph{[unit; \texttt{dax.py}]} \\
\texttt{test\_\allowbreak{}em\_\allowbreak{}recovers\_\allowbreak{}attractiveness\_\allowbreak{}and\_\allowbreak{}examination\_\allowbreak{}ranking} & Spearman above 0.8 for $\alpha$ and 0.75 for $\theta$; tab factors within 0.08 \emph{[unit; \texttt{dax.py}]} \\
\texttt{test\_\allowbreak{}em\_\allowbreak{}recovers\_\allowbreak{}examination\_\allowbreak{}up\_\allowbreak{}to\_\allowbreak{}the\_\allowbreak{}global\_\allowbreak{}scale} & Global scale removed with the true mean (test-only oracle); $\theta$ error below 0.05 \emph{[unit; \texttt{dax.py}]} \\
\texttt{test\_\allowbreak{}dax\_\allowbreak{}fit\_\allowbreak{}reads\_\allowbreak{}only\_\allowbreak{}pre\_\allowbreak{}cut\_\allowbreak{}rows\_\allowbreak{}and\_\allowbreak{}no\_\allowbreak{}labels} & Same parameters after adding post-cut rows or flipping every random label \emph{[unit; \texttt{dax.py}]} \\
\texttt{test\_\allowbreak{}additive\_\allowbreak{}arm\_\allowbreak{}tracks\_\allowbreak{}item\_\allowbreak{}appeal} & Additive appeal has Spearman above 0.5 with the truth \emph{[unit; \texttt{dax.py}]} \\
\texttt{test\_\allowbreak{}within\_\allowbreak{}tag\_\allowbreak{}control\_\allowbreak{}removes\_\allowbreak{}the\_\allowbreak{}item\_\allowbreak{}gain} & Within-tag permutation keeps under 30\% of the gain (stricter than gate B2) \emph{[unit; \texttt{dax.py}]} \\
\texttt{test\_\allowbreak{}permutations\_\allowbreak{}keep\_\allowbreak{}the\_\allowbreak{}tag\_\allowbreak{}or\_\allowbreak{}not} & Within-tag permutation keeps tags; all-item permutation is a bijection \emph{[unit; \texttt{dax.py}]} \\
\texttt{test\_\allowbreak{}channel\_\allowbreak{}interface\_\allowbreak{}and\_\allowbreak{}joint\_\allowbreak{}module\_\allowbreak{}shapes} & Channel shapes, interface and row standardisation \emph{[unit; \texttt{dax.py}]} \\
\texttt{test\_\allowbreak{}post\_\allowbreak{}hoc\_\allowbreak{}dax\_\allowbreak{}reproduces\_\allowbreak{}the\_\allowbreak{}probe\_\allowbreak{}split0} & Split-0 test value equals the probe value to $10^{-4}$ \emph{[unit; \texttt{probe\_\allowbreak{}pbm.json}]} \\
\multicolumn{2}{l}{\emph{Checks in the run logs and the code}} \\
Random log is not a training input & \texttt{training\_\allowbreak{}uses\_\allowbreak{}random\_\allowbreak{}log = False}; fits read only the pre-cut log \emph{[code; \texttt{data2.py}]} \\
Temporal boundary & Max pre-cut training time 1,650,556,797,718 ms $<T_0$ $\le$ min test time 1,650,595,332,831 ms; zero leak counts asserted \emph{[run; \texttt{dax\_\allowbreak{}lgn\_\allowbreak{}post.log}]} \\
Cached test lists equal the loader's & Test lists and users match the loader \emph{[run; same]} \\
Sequence builder cut-off & Sequence events kept only if time $<T_0$ and not on a pool video (timestamp, not file name) \emph{[code; \texttt{v12\_\allowbreak{}build\_\allowbreak{}seq.py}]} \\
Shuffled users, closed-form backbone & Full minus shuffled $+18.90$ [1.99] $\times10^{-3}$; shuffled minus no-lifelong $-1.02\times10^{-3}$, passes $\pm2\times10^{-3}$ \emph{[main; \texttt{stage1.txt}]} \\
Shuffled users, LEAF-LGN & Lifelong channels on shuffled tables: $-0.93$ [0.45] $\times10^{-3}$ \emph{[probe; \texttt{probe\_\allowbreak{}27k\_\allowbreak{}shuffled}]} \\
$\alpha$ permuted across all items & Gain $-0.27$ [0.36] $\times10^{-3}$; removes 101.8\% of the full gain (B2: at least 90\%) \emph{[main; \texttt{stageB.txt}]} \\
$\alpha$ permuted within tag & Gain $+6.28$ [1.10] $\times10^{-3}$; a drop of 58.4\% (B2: at least 30\%) \emph{[main; same]} \\
$\alpha$ tag mean; within tag; full & $+4.31$ [1.07]; $+3.29$ [0.99]; $+9.45$ [1.36] $\times10^{-3}$ \emph{[probe; \texttt{probe\_\allowbreak{}dax\_\allowbreak{}ablate}]} \\
$\alpha$ permuted across all items, weight & Selected weight 0 in 5 of 5 splits \emph{[probe; same]} \\
Full ablation arm reproduces the main run & Same configurations on all five splits; same test value \emph{[main; \texttt{dax\_\allowbreak{}ablate.json}]} \\
Stored LEAF-LGN reproduces & Post-hoc minus stored: $+0.73$ [0.52] $\times10^{-3}$, NB $p=0.660$; splits agree to 5 decimals \emph{[main; \texttt{stageB.txt}; \texttt{leaf\_\allowbreak{}scores}]} \\
User linkage, Pure to 27K & 1,186,059 random rows; identifiers identical for 0.9999 of joined events \emph{[check; \texttt{link\_\allowbreak{}users.txt}]} \\
Off-protocol calibration & Validation users' random rate added to LEAF-LGN: $+36.62$ [1.79] $\times10^{-3}$; post-cut labels; not a method \emph{[diagnostic; \texttt{probe\_\allowbreak{}items.txt}]} \\
\bottomrule
\end{tabular}
\end{adjustbox}
\end{table}

\subsection{Baselines}
\label{sec:protocol-baselines}

Table~\ref{tab:baselines} lists the baselines, their grids and their configuration counts. All baselines are trained on the same
pre-cut log with the same splits and the same selection rule. Every backbone with DAX gets the same post-hoc DAX module and the
same $\gamma$ grid as LEAF-LGN.

\begin{table}[t]
\centering
\caption{Baselines, search grids and configuration counts (the $n_{\mathrm{configs}}$ of the run logs, as in
Section~\ref{sec:protocol-splits}). ``Not run'' marks a row that has no result file.}
\label{tab:baselines}
\footnotesize
\begin{adjustbox}{max width=\textwidth}
\begin{tabular}{p{0.25\textwidth}p{0.45\textwidth}rl}
\toprule
Method & Search grid & Configurations & Source \\
\midrule
\raggedright LightGCN \citep{he2020lightgcn} & \raggedright lr $\{10^{-2},3\cdot10^{-3}\}$, wd $\{10^{-5},10^{-4}\}$; $d=64$, 2 layers, at most 40 epochs, patience 3 & 4 & \texttt{v10/learned\_kuairand\_t\_base.json} \\
\raggedright r-AdjNorm \citep{zhao2022radjnorm} & \raggedright LightGCN grid with normalisation exponent $r\in\{0.3,0.7\}$ & 8 & same \\
\raggedright iALS \citep{hu2008ials} (Stage 1 only) & \raggedright $\alpha\in\{3,10,30\}$, regularisation $\in\{1,10,100,1000\}$ & 12 & same \\
\raggedright QP family (Appendix A) & \raggedright DR-BPR-LGN-1, -p2 and -p with their own grids; per split the best of the three & 20 & \texttt{v10/learned\_kuairand\_t\_qp.json} \\
\raggedright QP family + DAX & \raggedright the 10 distinct (method, configuration) trainings chosen by any split of the QP run, with the DAX grid & 3,200 & \texttt{v11/dax\_qp.json} \\
\raggedright EASE \citep{steck2019ease} & \raggedright closed form, $\lambda\in\{1000,3000,10000\}$, $w\in\{0,0.25,0.5\}$, popularity $\beta\in\{0,0.25,0.5\}$ & 27 & \texttt{v11/dax\_ease.json} \\
\raggedright EASE + DAX & \raggedright EASE with the DAX grid & 8,640 & same \\
\raggedright LEAF-LGN (post-hoc backbone) & \raggedright 12 trainings (lr, wd, $w$), no DAX & 12 & \texttt{v11/dax\_lgn\_post.json} \\
\raggedright LEAF-LGN + DAX & \raggedright 12 trainings with the DAX grid and $\gamma$ & 3,840 & same \\
\raggedright DIN (deep interest network; \{TBD:citation key, not in the allowed bibliographies\}) & \raggedright 4 trainings (lr, wd); dim 16, hidden 16, MLP 32, batch 1024, at most 12 epochs, patience 2; last $L=200$ pre-cut events & 4 & \texttt{v12/din\_27k.json} \\
\raggedright DIN + DAX & \raggedright 4 trainings with the DAX grid and $\gamma$ & 1,280 & same \\
\raggedright SIM-hard \citep{pi2020sim} & \raggedright as DIN, with the $K=50$ most recent events sharing the tag or author of the candidate & 4 & \texttt{v12/sim\_1k.json} (1K) \\
\raggedright SIM-hard + DAX & \raggedright as DIN + DAX & 1,280 & same \\
\raggedright LEAF-tied (closed form, Stage 1 only) & \raggedright $\lambda$, $w$, $\beta$, $s$ and lr, wd & 3,456 & \texttt{v10/leaf\_tied.json} \\
\raggedright LightGCN + DAX & \raggedright not run & n/a & no result file \\
\bottomrule
\end{tabular}
\end{adjustbox}
\end{table}

The QP family is our own earlier method. It is not cited as a publication, and its full description is in Appendix A. The
QP family searches 20 configurations in total (4, 8 and 8 for its three members), and each split selects one of them on
validation, as described above.

The DIN and SIM baselines are lifelong-aware. Their sequence is the last $L=200$ pre-cut events of the user on videos outside the
pool, with the tag, author and music codes of each video and an engagement sign. DIN attends over the whole sequence. SIM-hard first
keeps the $K=50$ most recent events that share the tag or author of the candidate, and then attends over them. Both models also use
the identifier embeddings of users and items and are trained with the BPR loss on the pre-cut clicks, as LightGCN is. Their grids are
smaller than those of the graph backbones, with dimension 16 instead of 64, at most 12 epochs instead of 40, and patience 2 instead of 3.
The disclosure matters for the comparison: the DIN and SIM budgets are four trainings, and the results on them are conditional on that
budget. The SIM run on KuaiRand-27K did not finish. Its partial file holds three of the four trainings and no SIM result for that
setting is reported. SIM is therefore reported on KuaiRand-1K only.
% source: experiments/run_din.py (grid, L, sim_k); drup/din.py (score, sparse embeddings); experiments/v12_build_seq.py; results/v12/sim_27k.json.partial.json

The NAVIP method \citep{kim2022navip} is not re-run. It changes the neighbour aggregation of a graph network and does not use
lifelong histories, and no run of it is reported. TWIN \citep{chang2023twin} and the soft-search variant of SIM are not run, and the
comparisons in the paper do not include them.

\subsection{Ablation design}
\label{sec:protocol-ablation}

The ablations of DAX are a cut-down version of the main search. Each ablation arm is a separate method and is selected per split in the
same way as the full arm. The arms are run on the backbone settings that the main run chose: each split reads the epochs within three
epochs of its chosen epoch, and the trainings are the two settings that the splits chose. The arms use five values of $\gamma$,
$\{0.1,0.25,0.5,0.7,1.0\}$, instead of the ten of the main run. The arms are the three DAX switches (no tab factor, no lifelong prior,
no tag prior; each with 16 configurations), the additive MRP-style logistic with tab effects (32 configurations, matched to the 32 of the
full grid), and the two permutations of the attractiveness: within tag (the tag mean is kept) and across all items. Each permutation arm
is run with three permutation seeds.

Because the trainings and epoch windows are restricted, the base of the ablation run is its own LEAF-LGN, which is selected inside the
window and has a lower value than the main-run LEAF-LGN. The ablation gains are paired differences against that base, so they are not the
same quantity as the main-table DAX gain against the main-run backbone. The full arm of the ablation run selects the same configurations on all five
splits as the main run and reproduces its test value exactly (Table~\ref{tab:controls}).

The gate B2 conditions are fixed in \texttt{paper/v7/DESIGN.md} (Section 7): the within-tag permutation must lower the gain by at least
30\%, the across-item permutation must remove at least 90\% of it, and the additive arm must stay at least $4\times10^{-3}$ below the full
DAX arm. The stored analysis reports the raw paired p-values of the ablation arms and does not adjust them. The Holm adjustment within the
ablation family is therefore not part of the reported results, and the gate decisions do not use p-values.
% source: experiments/run_dax.py (--ablate, --epoch_window, --gammas, --perm_seeds); experiments/final_v11.py (B2 thresholds and checks)

\subsection{Kinds of number}
\label{sec:protocol-kinds}

Three kinds of number are used, and they are not mixed within one table without a label. A main number is a validation-selected result
of the main setting, from the files in \texttt{results/v10/} and \texttt{results/v11/}, and of the v12 deep baselines. A probe number is a
post-hoc computation on stored scores, from \texttt{results/v11\_probes/} and \texttt{results/v9\_probes/}. A diagnostic number uses
random-exposure labels of the validation users and is never a method or a selection input. A number that is not in a result file is marked
as TBD in the text, and the table entry names the file that would supply it.
